\section{Related Work}
\label{sec:related}

\paragraph{Spatially adaptive reconstruction.}
ClassSR routes super-resolution patches to networks of different
capacity~\cite{kong2021classsr}; APE instead permits patches to exit a shared
network at different depths~\cite{wang2022ape}. These are direct precedents
for spatial capacity selection, so patch routing itself is not our novelty
claim. We examine the additional constraints of a coded representation:
which decision information is available at each endpoint, what must be
signalled, and how a local assignment affects final reconstruction quality.

\paragraph{Adaptive complexity in compression.}
Slimmable compressive autoencoders provide multiple model
widths~\cite{yang2021slimcae}. AdaNIC combines blockwise transform capacity
with a routing agent~\cite{tao2023adanic}, while cgSlimDecoder selects
decoder channel widths under complexity constraints~\cite{hu2023cgslim}.
These methods already establish content- or complexity-dependent computation
within learned compression. Our comparison focuses on depth allocation in
the synthesis path of a fixed coded representation and on its incremental
benefit over blind mixtures. A matched reproduction of these alternatives
remains necessary for a claim of superiority over adaptive codec designs.

\paragraph{Efficient synthesis and practical execution.}
Shallow decoders investigate how computation can be shifted toward the
encoder~\cite{yang2023shallow}. DCVC-RT distinguishes arithmetic cost from
operational overhead, including memory traffic and function
calls~\cite{jia2025dcvcrt}; DCVC-UF further develops a throughput-oriented
video codec~\cite{li2026dcvcuf}. These results motivate reporting actual
execution in addition to multiply--accumulates. Our scope is the supplied
intra model, not DCVC-UF's inter-frame or chunk-coding performance.

\paragraph{Distortion and perception.}
Generative compression explicitly studies reconstruction fidelity and
perceptual realism~\cite{mentzer2020,muckley2023illm}. A small PSNR change
therefore cannot establish perceptual equivalence. We use archived
non-DCVC experiments only to examine this limitation and the dependence on
decoder architecture; they are not evidence of a perceptual guarantee for
FLEX.
